% =============================================================================
\section{Tarefa 5 --- Fusão das duas estimativas}
\label{sec:fusao}
% =============================================================================

As duas estimativas de inclinação têm defeitos complementares: a do
acelerómetro não deriva mas é ruidosa e falha quando o sensor acelera; a do
giroscópio é suave e imune à aceleração, mas deriva.

\subsection{Covariance Intersection}

Sejam $\hat{\vect x}_a$ e $\hat{\vect x}_g$ as duas estimativas de
$(\phi;\,\theta)$, com covariâncias $P_a$ e $P_g$. A
\emph{Covariance Intersection}~\cite{julier1997} combina-as com um peso
$w \in [0,1]$:
\begin{equation}
P^{-1} = w\,P_a^{-1} + (1-w)\,P_g^{-1}, \qquad
\hat{\vect x} = P\bigl(w\,P_a^{-1}\hat{\vect x}_a + (1-w)\,P_g^{-1}\hat{\vect x}_g\bigr).
\label{eq:ci}
\end{equation}
Com $w=1$ fica só o acelerómetro e com $w=0$ só o giroscópio. Ao contrário
de um filtro de Kalman, a CI não assume que os erros das duas estimativas são
independentes, e o resultado nunca é mais confiante do que os dados
justificam. Aqui isso é relevante: a estimativa fundida de cada instante é o
ponto de partida da integração seguinte, e os erros do acelerómetro (desvio
residual, acelerações que duram várias amostras) estão correlacionados no
tempo.

\todo{como se obtêm $P_a$ (ruído do acelerómetro propagado
por~\eqref{eq:incl}; aumentada quando o módulo se afasta de $1\,g$) e $P_g$
(cresce a cada passo com $\sigma_g$); se $w$ é fixo ou escolhido a cada
instante para minimizar o traço de $P$.}

\subsection{Resultados}

\begin{figure}[H]
\centering
\fig{Photos/fusao_tempo.png}{\textwidth}
\caption{\todo{um varrimento de inclinação (ensaios D e D2): ângulo imposto
pelo robô, estimativa do acelerómetro, do giroscópio e da fusão; por baixo, o
erro de cada uma.}}
\label{fig:fusao}
\end{figure}

\begin{figure}[H]
\centering
\fig{Photos/fusao_w.png}{0.7\textwidth}
\caption{\todo{erro quadrático médio da fusão em função de $w$, com o mínimo
assinalado; estimativas no tempo para três valores de $w$.}}
\label{fig:ciw}
\end{figure}

\begin{table}[H]
\centering
\caption{Erro de inclinação de cada método contra a referência do robô.
\todo{preencher.}}
\label{tab:fusao}
\small
\begin{tabular}{lcccc}
\toprule
\rowcolor{tabcab}
 & \multicolumn{2}{c}{Ensaio D} & \multicolumn{2}{c}{Ensaio D2 (com perturbação)} \\
\rowcolor{tabcab}
Método & RMS & máx. & RMS & máx. \\
\midrule
Acelerómetro (Tarefa 3) & & & & \\
Giroscópio (Tarefa 4)   & & & & \\
CI, melhor $w$          & & & & \\
Filtro de Kalman        & & & & \\
\bottomrule
\end{tabular}
\end{table}

\subsection{Comparação com um filtro de Kalman}

\todo{opcional, a confirmar com o docente. As Tarefas 3 e 4 são as duas
metades de um filtro de Kalman~\cite{kalman1960}: o giroscópio faz a previsão
e o acelerómetro a correção. Estado: orientação e desvio do giroscópio.
Comparar com a CI nos mesmos dados (Tabela~\ref{tab:fusao}) e discutir a
confiança de cada um (covariância estimada contra erro real).}
